\section{Independent Follow-up with Available Local Artifacts}
\label{app:followup}

Company policy prevents transfer of the ImageNet checkpoints. We therefore
separate two feasible follow-ups: architecture timing using random weights,
and prediction tests using existing, locally available CIFAR-100 weights.
Neither experiment independently reproduces the company ImageNet scores.

\subsection{A common A100 inference harness}

We compare dense DeiT-S/8, ToMe, and the canonical sparse-backend \mn{} on
an A100-SXM4-80GB using PyTorch 2.6.0+cu124, timm 0.9.11 and
FlashAttention 2.7.4.post1. All cells use batch 64, float32 parameters and
inputs, and fp16 autocast. The compressed models retain exactly 392 patches
at 224 pixels and 1152 at 384 pixels, excluding CLS. A diagnostic forward
records block outputs; all hooks are removed before timing. ToMe source
tracing is disabled. Proportional attention is explicitly controlled, with
its enabled setting matching the recorded accuracy protocol.

Each cell has ten warm-up iterations and thirty CUDA-event samples in each
of three rounds, reversing model order in the middle round. A selected-GPU
process guard requires an empty device before setup and only our process
during the experiment. Models are allocated and released sequentially.
Table~\ref{tab:followup-latency} reports the median of three round medians.
It measures model execution, excluding data loading, backward and updates.

\begin{table}[h]
\centering\small
\begin{tabular}{lrr}
\toprule
Model & 224 px (ms) & 384 px (ms) \\
\midrule
Dense DeiT-S/8 & 30.09 & 109.43 \\
ToMe, proportional attention on & 34.42 & 133.36 \\
ToMe, proportional attention off & 29.71 & 96.05 \\
MergeNet $R=3$ & 49.16 & 144.97 \\
MergeNet $R=3$, BNHD & 46.56 & 138.11 \\
MergeNet $R=5.143$ & -- & 218.25 \\
MergeNet $R=5.143$, BNHD & -- & 211.40 \\
\bottomrule
\end{tabular}

\caption{Independent synthetic A100 inference, batch 64. All weights are
randomly initialized; these are not trained-checkpoint frontier points.
BNHD changes the local-attention tensor layout. The 384-pixel radius-5.143
model has the same final patch budget as radius 3.}
\label{tab:followup-latency}
\end{table}

On this local hardware/software stack, the original \mn{} is slower than
dense DeiT at both resolutions. The BNHD change reduces complete-model
radius-3 time by about 5\%, considerably less than the isolated-module
improvement. The proportional-attention setting materially affects ToMe
latency. Thus neither the H20 crossover nor a ToMe timing measured with
proportional attention disabled should be treated as accelerator-independent
evidence. The A100 comparison changes the hardware and software stack
relative to the H20 experiment; it does not isolate a hardware-only cause.
No historical accuracy point is attached to these random-weight timings.

\subsection{Layout validation under a complete training graph}

The layout comparison uses identical weights and images, restored CPU and
CUDA random states, fresh per-sample routing, and zero dropout at both
224 and 384 pixels. With the dense routing reference backend, outputs,
input gradients and all parameter gradients are bitwise equal. With the
sparse backend and an unscaled squared-logit loss, the strict gradient
comparison fails; an unchanged-model repeat also fails, so these runs do
not isolate an error introduced by BNHD. Their input-gradient relative
errors are approximately 4.8\% and 16.3\%, respectively.

We retain these failed tests and add a fixed loss scale of 1024, unscaling
all gradients before comparison. The sparse-backend tests then pass the
predefined relative-$L_2$ limits of 0.002 for outputs and 0.003 for gradients.
Maximum parameter-gradient relative errors are below 0.0016; unchanged-model
repeat controls are below 0.0018. This diagnoses sensitivity of the unscaled
fp16 test and supports bounded layout parity. It does not establish long-run
training equivalence, eliminate sparse-kernel nondeterminism, or validate
arbitrary dropout and precision settings.
